% !TEX root = ../my-thesis.tex
%
\pdfbookmark[0]{Abstract}{Abstract}
\addchap*{Abstract}
\label{sec:abstract}
Diese Arbeit vergleicht die fünf Workflow-Management-Systeme
Nextflow, StreamFlow, Merlin, Pegasus und AiiDA anhand zweier
Workflow-Muster, einer Pipeline und eines Scatter-Gather-Musters,
die in allen Systemen mit derselben Rechenlogik umgesetzt werden.
Betrachtet werden die Modellierung der Workflows, das
Laufzeitverhalten anhand von Makespan, Overhead und
Koordinationsgrößen sowie der Betrieb in einer HPC-Umgebung.
Zunächst werden alle fünf Systeme unter kontrollierten Bedingungen
auf einem einzelnen Rechner verglichen. Anschließend wird der
Einsatz auf dem HPC-System MOGON unter Slurm betrachtet, auf dem
vier der fünf Systeme betrieben werden konnten.
Beide Workflow-Muster ließen sich in allen fünf Systemen umsetzen.
Die wesentlichen Unterschiede zeigen sich in der Art der
Workflow-Beschreibung, im Laufzeitverhalten und bei der Koordination
der Schritte sowie im Betriebsmodell. Beim Laufzeitverhalten liegen StreamFlow, Nextflow und Merlin beim Overhead in einem ähnlichen Bereich, während AiiDA höhere Werte
aufweist und Pegasus in der lokalen Untersuchung deutlich darüber
liegt. Mit zunehmender Rechenlast verliert der Overhead gegenüber
der eigentlichen Rechenzeit an Bedeutung. Auf MOGON zeigte sich
zudem, dass durch die Clusterumgebung verursachte Zeitanteile die
Unterschiede zwischen den Systemen überlagern können. Insgesamt
unterscheiden sich die Systeme sowohl in der Modellierung und
Ausführung der Workflows als auch in ihrem Betrieb in der
HPC-Umgebung.